\chapter{Bevezetés}
\label{ch:intro}

A keresőalgoritumusok meghatározó szerepet játszanak az informatikában és azon belül a mesterséges intelligencia területén. Ezek az algoritmusok elengedhetetlenek az útkeresési feladatokban, a játékok mint a sakk elemzésében, az automatikus következtetésben is. Bár ezeknek az algoritmusoknak az elméleti tulajdonságai, mint például a teljesség, optimalitás és a számítási komlexitás már jól ismertek, maguk az algoritmusoknak a megértésének elsajátításához nem találunk sok eszközt. 
Azonban az felfedezett állapottér ábrázolásával sokkal világosabbá válhat, hogy adott kereső algoritmusok, milyen szabályok szerint hoznak döntéseket, a globális munkaterét ábrázolva, megfigyelhető, hogy milyen megszerzett tudással tud dolgozni az algoritmus. Kielemezhető, hogy milyen szerepet játszanak a jó heurisztikus függvények, és ez által milyen eredményeket produkálnak különböző problémákon.

Ez a dolgozat egy olyan interaktív alkalmazást mutat be, amely keresőalgoritmusokat alkalmaz a Rush Hour logikai játékra~\cite{rushhour_puzzle}, és a feltárt állapotgráfot háromdimenziós térben vizualizálja. A Rush Hour egy olyan rejtvény, amelyben a játékosnak egy $6 \times 6$-os táblán kell járműveket csúsztatnia, hogy utat törjön a fő (piros) autónak a kijárat felé. A játék szabályai nagyon egyszerűek, viszont a feladványok komplexitása a viszonylag könnyűtől a meglepően bonyolult nehézségig terjed, ezért ideális választásnak bizonyult az elemzésre kerülő problémaként.  

A projekt két oldalról közelíti meg a problémát:

\begin{enumerate}
    \item \textbf{Keresőalgoritmusok implementálása és vezérlése.} Három keresőalgoritmus került implementálásra: a hegymászó keresés általános verziója - a tabu keresés, visszalépéses (backtracking) keresés és A$^C$ (A és A$^*$ speciális változata) gráfkeresés, amelyek mindegyike konfigurálható a négy heurisztikus függvény egyikével és egy opcionális véletlenszerűségi beállítással. A felhasználó manuálisan is játszhat a rejtvénnyel, és minden felfedezett állapot megjelenik az állapotgráfban és minden szomszédos állapot össze lesz kötve.
    \item \textbf{Háromdimenziós állapotgráf-vizualizáció.} A felfedezett állapotok kör alakú csúcsként jelennek meg a három dimenziós térben, amelyeket az érvényes lépéseket reprezentáló élek kötnek össze. Egy erő-alapú elrendezési (force-directed layout) fizikai szimuláció mozgatja a csúcsokat úgy, hogy a szomszédos állapotok viszonylag közel maradjanak egymáshoz, míg a nem szomszédosak távolodjanak, ehhez az élek mentén rugóerőket, a csúcspárok között pedig a távolság négyzetének reciprokával arányos taszítóerőt alkalmazunk. A taszítási számítások $O(V^2)$ költségének kezelése érdekében a Barnes–Hut optimalizációs algoritmust~\cite{barneshut1986} és egy nyolcasfa (octree) térbeli adatszerkezetet használunk, ami az időbeli komplexitást képkockánként $O(V \log V)$-re csökkenti. (V a felfedezett állapotok azaz a megjelenített csúcsok száma)
\end{enumerate}

Az alkalmazás Godot Engine 4.6~\cite{godot_engine} használatával, C\# és .NET~\cite{dotnet10} nyelven készült, és a Modell–Nézet (Model–View) arkitektúrát alkalmazza. A modell réteg tiszta C\# konstrukciókkal enkapszulálja a rejtvény logikáját, a keresőalgoritmusokat és a heurisztikus függvényeket, míg a nézet réteg a Godot jelenetrendszerét használja a megjelenítéshez, a fizikai szimulációhoz és a felhasználói interakcióhoz.

